\chapter{Sepsis}

\section*{Definition}
Sepsis is life-threatening organ dysfunction caused by a dysregulated response to infection. It is a medical emergency. Septic shock is a severe form of sepsis with persistent circulatory and cellular abnormalities associated with a particularly high risk of death.

\section*{When to See a Doctor}
Call emergency services immediately for suspected sepsis, especially when infection is accompanied by confusion, extreme drowsiness, difficulty breathing, blue or mottled skin, fainting, severe weakness, very fast breathing or heartbeat, cold clammy skin, little or no urine, or rapidly worsening illness. Fever may be absent, particularly in older or immunocompromised people. Do not wait for all signs to appear.

\section*{How It Develops}
An infection triggers an inflammatory and immune response that can become harmful to the body's own tissues. Widespread changes in blood-vessel tone, capillary permeability, clotting, and cellular metabolism impair circulation and oxygen use. The resulting organ dysfunction may affect the brain, lungs, kidneys, liver, heart, or blood-clotting system.

\section*{Causes}
\begin{itemize}
  \item Bacterial, viral, fungal, or parasitic infections; pneumonia, urinary or kidney infection, abdominal infection, skin or soft-tissue infection, and meningitis are common sources.
  \item Infection of a surgical site, infected wound, abscess, intravenous catheter, prosthetic device, or other healthcare-associated source.
  \item Risk is higher in infants, older adults, pregnant or postpartum people, people with weakened immunity, chronic disease, recent surgery, or indwelling devices.
  \item The source may not be immediately apparent, especially in older adults or people taking immune-suppressing medicines.
\end{itemize}

\section*{Related Symptoms}
\begin{itemize}
  \item Fever or unusually low temperature, chills, sweating, rapid breathing, rapid heart rate, severe pain, or profound malaise.
  \item Confusion, reduced alertness, dizziness, low blood pressure, mottled skin, reduced urine, nausea, vomiting, or diarrhea.
  \item Cough, shortness of breath, painful urination, flank pain, abdominal pain, wound drainage, rash, or another symptom indicating the infection source.
\end{itemize}

\section*{Medical Evaluation}
\begin{itemize}
  \item Immediate assessment includes airway, breathing, circulation, mental status, temperature, blood pressure, respiratory rate, oxygen saturation, urine output, and peripheral perfusion.
  \item Blood tests commonly include lactate, complete blood count, electrolytes, kidney and liver function, glucose, coagulation studies, and blood cultures. Cultures should be obtained promptly when feasible without delaying antibiotics.
  \item Urine testing, chest imaging, ultrasound, CT, or other targeted studies are selected to identify the infection source and complications.
  \item Clinicians assess for organ dysfunction, including hypoxemia, hypotension, acute kidney injury, altered mental status, low platelets, abnormal clotting, or rising lactate.
\end{itemize}

\section*{Treatment Options}
\begin{itemize}
  \item Begin prompt intravenous antimicrobial treatment for suspected sepsis after appropriate cultures when this does not cause a clinically important delay.
  \item Give intravenous crystalloid for sepsis-related hypoperfusion, with frequent reassessment to avoid fluid overload; use vasopressors, usually norepinephrine, when hypotension persists.
  \item Control the source through drainage of an abscess, removal of an infected device, surgery, or another appropriate procedure.
  \item Provide oxygen and organ support as needed, including ventilatory support, renal replacement therapy, blood products, glucose management, and intensive-care monitoring.
  \item Narrow or stop antimicrobials when microbiology and clinical response identify a safer targeted approach.
\end{itemize}

\section*{Self-Care and Home Management}
Sepsis cannot be safely treated at home. Call emergency services, describe the suspected infection and warning signs, and bring a medication list and relevant medical history. Do not give food, drink, or oral medicines to a person who is confused, vomiting, unconscious, or unable to swallow safely. Do not delay transport while trying to reduce a fever.

\section*{References}
\begin{thebibliography}{99}
\bibitem{sepsis:ref1} Evans L, Rhodes A, Alhazzani W, et al. Surviving Sepsis Campaign: international guidelines for management of sepsis and septic shock 2021. \textit{Intensive Care Med}. 2021;47:1181--1247.
\bibitem{sepsis:ref2} Singer M, Deutschman CS, Seymour CW, et al. The Third International Consensus Definitions for Sepsis and Septic Shock (Sepsis-3). \textit{JAMA}. 2016;315:801--810.
\bibitem{sepsis:ref3} World Health Organization. Sepsis. WHO; updated 2024.
\end{thebibliography}
